\cotefile{156-4}
\scanpage{10}
\mspage{9}

Ainsi, on est \unsure{arrivé|amené} à la structure suivante

\textcircled{1} $\unsure{\mathcal{F}|F}_{\unsure{el|\text{él}}}$ \unsure{ens|ens.} des ``figures élémentaires''

\textcircled{2} relation d'ordre sur $\mathcal{F}_{el}$, $\leq$ (inclusion des figures)

\textcircled{3} relation symétrique réflexive $R$ \deleted{\unsure{$F,G$|$F\,R\,G$}} ``$F$ et $G$ \uline{\unsure{compatibles}}''.\ \deleted{\illegible{1}} \deleted{\unsure{Parties|Partie}}

\margin{\uline{Configuration} (\unsure{élémentaire}) partie de $\mathcal{F}_{\text{él}}$ \unsure{formée|fermée} d'él. deux \unsure{à|:} deux \unsure{compatibles|comp.tibles}}

% slanted
\margin{deux configurations sont dites compatibles si leur réunion l'est.}

\textcircled{4} \deleted{Antifiltre sur $\mathcal{F}_{el}$, des parties} Parmi les configurations, \deleted{\unsure{descripta|description|decripta}} \inserted{donnée de l'ens.} \deleted{\unsure{loc. finies}} \deleted{\unsure{prop.}} \deleted{\illegible{1}} de celles qui sont ``loc. finies'', avec \deleted{Avec} propriétés d'antifiltre \unsure{;|,} $\mathrm{Conf}_{\mathrm{l.f.}}\unsure{\subset|\in}\mathcal{P}(\mathcal{F}_{el})$

a) Toute configuration réduite \unsure{à|:} une seule figure él. est loc. finie

% slanted
\margin{Cas important (cas ``\uline{compact}'') : les configurations loc. finies sont les configurations finies.}

b) Si $\Phi$ est une configuration, $\Phi$ loc. finie ssi la \deleted{\unsure{réduite|réduit}} partie \unsure{fermée|fermée.} (pour la rel. d'ordre $\leq$) l'est

c) Si $\Phi$ est une configuration, \deleted{$\Phi=\bigcup_{i}$} loc. finie, alors tte $\unsure{\Psi|\psi}\subset\Phi$ est loc. finie

d) Si $\unsure{\Phi|\Psi}$ est une configuration, $\Phi=\bigcup_{i}\Phi_i$ (réunion finie) si les $\Phi_i$ sont loc. \unsure{finis|finies}, alors $\Phi$ aussi.
